% Part: first-order-logic
% Chapter: syntax-and-semantics
% Section: formation-sequences

\documentclass[../../../include/open-logic-section]{subfiles}

\begin{document}

\olfileid{fol}{syn}{fseq}

\olsection{నిర్మాణ క్రమాలు}

ఆగమనాత్మక నిర్వచనం ద్వారా \tetoken{సూత్రాలను}{!!{formula}s} నిర్వచించడం, దానికి పూరకంగా
ఆగమనం ద్వారా \tetoken{సూత్రాలు}{!!{formula}s} ధర్మాలను నిరూపించడం సొగసైన, సమర్థమైన
విధానం. అయితే \tetoken{సూత్రాలు}{!!{formula}s} నిర్మాణాన్ని కింది నుంచి పైకి, మెట్టుమెట్టుగా
పరిశీలించడం కూడా ఉపయోగకరం కావచ్చు. ఇక్కడ \emph{నిర్మాణ క్రమం} అనే
భావనతో అదే చేస్తాం.
%
పదాలు, \tetoken{సూత్రాలను}{!!{formula}s} వాటి ఆగమనాత్మక నిర్వచనాలను ప్రస్తావించకుండానే
ఈ విధంగా ఎలా ప్రవేశపెట్టవచ్చో చూపడానికి, ముందుగా ఏదో భాష~$\Lang L$
నుంచి తీసుకున్న సంకేతాల యాదృచ్ఛిక సంకేతమాల అనే భావనను ప్రవేశపెడతాం.

\begin{defn}[సంకేతమాలలు]
\ollabel{defn:string}
$\Lang L$ ఒక మొదటిస్థాయి విధేయ భాష అనుకుందాం. \emph{$\Lang
L$-సంకేతమాల} అంటే $\Lang L$లోని సంకేతాల పరిమిత క్రమం. భాష~$\Lang L$
ఏదో సందర్భం నుంచి స్పష్టంగా ఉన్నప్పుడు, $\Lang L$-సంకేతమాలను తరచుగా
కేవలం \emph{సంకేతమాల} అంటాం.
\end{defn}

\begin{ex}
ఏ మొదటిస్థాయి విధేయ భాష~$\Lang L$కైనా, అన్ని $\Lang L$-\tetoken{సూత్రాలు}{!!{formula}s}
$\Lang L$-సంకేతమాలలు; కానీ తిరుగుదిశ నిజం కాదు. ఉదాహరణకు,
\[)(\Obj v_0\lif\lexists\] అనేది $\Lang L$-సంకేతమాల, కానీ
$\Lang L$-\tetoken{సూత్రం}{!!{formula}} కాదు.
\end{ex}

\begin{defn}[పదాల నిర్మాణ క్రమాలు]
\ollabel{defn:fseq-trm}
$\Lang L$-సంకేతమాలల పరిమిత క్రమం $\tuple{t_0,\dotsc,t_n}$, పదం~$t$కు
\emph{నిర్మాణ క్రమం} కావడం అంటే $t \ident t_n$, ఇంకా ప్రతి $i \leq n$కు
$t_i$ \tetoken{ఒక చరం}{!!a{variable}} లేదా \tetoken{ఒక స్థిరసంకేతం}{!!a{constant}} అయి ఉండాలి; లేదా $\Lang L$లో
$k$-స్థాన \tetoken{ప్రమేయ సంకేతం}{!!{function}}~$f$ ఉండి, $t_i \ident
f(t_{m_1},\dotsc,t_{m_k})$ అయ్యే $m_1,\dotsc,m_k < i$లు ఉండాలి.
\sourcecorrection{OLTEFOLSYN-003}{ఆంగ్ల మూలం $k$-స్థాన ప్రమేయానికి
$m_0,\dotsc,m_k$లతో $k+1$ ఆర్గ్యుమెంట్లను ఇచ్చింది. ఇదే అధ్యాయంలోని
$k$-స్థాన నిర్వచనానికి సరిపడేలా సూచికలను $m_1,\dotsc,m_k$గా చేశాం.}
భేదం చెప్పాల్సినప్పుడు, పదాల నిర్మాణ క్రమాలను \emph{పద నిర్మాణ
క్రమాలు} అంటాం.
\end{defn}

\begin{ex}
క్రమం
\[
    \tuple{\Obj c_0, \Obj v_0, \Atom{\Obj f^2_0}{\Obj c_0, \Obj v_0}, \Atom{\Obj f^1_0}{\Atom{\Obj f^2_0}{\Obj c_0, \Obj v_0}}}
\]
అనేది $\Atom{\Obj f^1_0}{\Atom{\Obj
f^2_0}{\Obj c_0, \Obj v_0}}$ పదానికి నిర్మాణ క్రమం; కిందిదీ అలాంటిదే:
\[
    \tuple{\Obj v_0, \Obj c_0, \Atom{\Obj f^2_0}{\Obj c_0, \Obj v_0}, \Atom{\Obj f^1_0}{\Atom{\Obj f^2_0}{\Obj c_0, \Obj v_0}}}.
\]
\end{ex}

\begin{defn}[సూత్రాల నిర్మాణ క్రమాలు]
\ollabel{defn:fseq-frm}
$\Lang L$-సంకేతమాలల పరిమిత క్రమం $\tuple{!A_0,\dotsc,!A_n}$,
$!A$కు \emph{నిర్మాణ క్రమం} కావడం అంటే $!A \ident !A_n$, ఇంకా
ప్రతి $i \leq n$కు $!A_i$ ఒక పరమాణు \tetoken{సూత్రం}{!!{formula}} అయి ఉండాలి; లేదా
కింది వాటిలో ఒకటి నిజమయ్యేలా $j,k < i$లు, \tetoken{ఒక చరం}{!!a{variable}}~$x$ ఉండాలి:
\begin{enumerate}
    \tagitem{prvNot}{$!A_i \ident \lnot !A_j$.}{}%
    \tagitem{prvAnd}{$!A_i \ident (!A_j \land !A_k)$.}{}%
    \tagitem{prvOr}{$!A_i \ident (!A_j \lor !A_k)$.}{}%
    \tagitem{prvIf}{$!A_i \ident (!A_j \lif !A_k)$.}{}%
    \tagitem{prvIff}{$!A_i \ident (!A_j \liff !A_k)$.}{}%
    \tagitem{prvAll}{$!A_i \ident \lforall[x][!A_j]$.}{}%
    \tagitem{prvEx}{$!A_i \ident \lexists[x][!A_j]$.}{}%
\end{enumerate}
భేదం చెప్పాల్సినప్పుడు సూత్రాల నిర్మాణ క్రమాలను \emph{సూత్ర నిర్మాణ
క్రమాలు} అంటాం.
\end{defn}

\begin{ex}
\[
\tuple{
    \Atom{\Obj A^1_0}{\Obj v_0},
    \Atom{\Obj A^1_1}{\Obj c_1},
    (\Atom{\Obj A^1_1}{\Obj c_1} \land \Atom{\Obj A^1_0}{\Obj v_0}),
    \lexists[\Obj v_0][(\Atom{\Obj A^1_1}{\Obj c_1} \land \Atom{\Obj A^1_0}{\Obj v_0})]
}
\]
అనేది $\lexists[\Obj v_0][(\Atom{\Obj A^1_1}{\Obj c_1} \land \Atom{\Obj
A^1_0}{\Obj v_0})]$కు నిర్మాణ క్రమం; కిందిదీ అలాంటిదే:
\begin{multline*}
\tuple{
    \Atom{\Obj A^1_0}{\Obj v_0},
    \Atom{\Obj A^1_1}{\Obj c_1},
    (\Atom{\Obj A^1_1}{\Obj c_1} \land \Atom{\Obj A^1_0}{\Obj v_0}),
    \Atom{\Obj A^1_1}{\Obj c_1},\\
    \lforall[\Obj v_1][\Atom{\Obj A^1_0}{\Obj v_0}],
    \lexists[\Obj v_0][(\Atom{\Obj A^1_1}{\Obj c_1} \land \Atom{\Obj A^1_0}{\Obj v_0})]
}.
\end{multline*}
%
రెండవ ఉదాహరణలో కనిపించినట్లుగా, నిర్మాణ క్రమాల్లో ``వ్యర్థ'' భాగాలు
ఉండవచ్చు: అవి అనవసరమైన లేదా నిర్మాణానికి తోడ్పడని \tetoken{సూత్రాలు}{!!{formula}s}.
\end{ex}

\begin{prop}\ollabel{prop:formed}
$\Frm[L]$లోని ప్రతి \tetoken{సూత్రం}{!!{formula}}~$!A$కు ఒక నిర్మాణ క్రమం ఉంటుంది.
\end{prop}

\begin{proof}
$!A$ పరమాణువు అనుకుందాం. అప్పుడు $\tuple{!A}$ అనే క్రమం $!A$కు
నిర్మాణ క్రమం.
%
ఇప్పుడు $!B$,~$!C$లకు వరుసగా $\tuple{!B_0,\dotsc,!B_n}$,
$\tuple{!C_0,\dotsc,!C_m}$ అనే నిర్మాణ క్రమాలు ఉన్నాయనుకుందాం.
%
\begin{enumerate}
  \tagitem{prvNot}{$!A \ident \lnot !B$ అయితే,
      $\tuple{!B_0,\dotsc,!B_n,\lnot !B_n}$ అనేది $!A$కు నిర్మాణ
      క్రమం.}{}
  \tagitem{prvAnd}{$!A \ident (!B \land !C)$ అయితే,
      $\tuple{!B_0,\dotsc,!B_n,!C_0,\dotsc,!C_m,(!B_n \land !C_m)}$
      అనేది $!A$కు నిర్మాణ క్రమం.}{}
  \tagitem{prvOr}{$!A \ident (!B \lor !C)$ అయితే,
      $\tuple{!B_0,\dotsc,!B_n,!C_0,\dotsc,!C_m,(!B_n \lor !C_m)}$
      అనేది $!A$కు నిర్మాణ క్రమం.}{}
  \tagitem{prvIf}{$!A \ident (!B \lif !C)$ అయితే,
      $\tuple{!B_0,\dotsc,!B_n,!C_0,\dotsc,!C_m,(!B_n \lif !C_m)}$
      అనేది $!A$కు నిర్మాణ క్రమం.}{}
  \tagitem{prvIff}{$!A \ident (!B \liff !C)$ అయితే,
      $\tuple{!B_0,\dotsc,!B_n,!C_0,\dotsc,!C_m,(!B_n \liff !C_m)}$
      అనేది $!A$కు నిర్మాణ క్రమం.}{}
  \tagitem{prvAll}{$!A \ident \lforall[x][!B]$ అయితే,
      $\tuple{!B_0,\dotsc,!B_n,\lforall[x][!B_n]}$ అనేది $!A$కు
      నిర్మాణ క్రమం.}{}
  \tagitem{prvEx}{$!A \ident \lexists[x][!B]$ అయితే,
      $\tuple{!B_0,\dotsc,!B_n,\lexists[x][!B_n]}$ అనేది $!A$కు
      నిర్మాణ క్రమం.}{}
\end{enumerate}
\tetoken{సూత్రాలపై}{!!{formula}s} ఆగమన సూత్రం ప్రకారం ప్రతి \tetoken{సూత్రానికి}{!!{formula}} ఒక నిర్మాణ
క్రమం ఉంటుంది.
\end{proof}

తిరుగుదిశను కూడా నిరూపించవచ్చు. సూత్రాలను నిర్వచించే మన రెండు
పద్ధతులు తుల్యమైనవని---అవి ఒకే ఫలితాలను ఇస్తాయని---ఇది చూపుతుంది
కనుక ముఖ్యం. నిర్మాణ క్రమాల పొడవుపై సాధారణ ఆగమనాన్ని ఉపయోగించి
సూత్రాల గురించిన సిద్ధాంతాలను నిరూపించవచ్చని కూడా దీని అర్థం.

\begin{lem}
\ollabel{lem:fseq-init}
$\tuple{!A_0,\dotsc,!A_n}$ అనేది $!A_n$కు నిర్మాణ క్రమమనీ,
$k \leq n$ అనీ అనుకుందాం. అప్పుడు $\tuple{!A_0,\dotsc,!A_k}$
అనేది $!A_k$కు నిర్మాణ క్రమం.
\end{lem}

\begin{proof}
సాధన.
\end{proof}

\begin{prob}
\olref[fol][syn][fseq]{lem:fseq-init}ను నిరూపించండి.
\end{prob}

\begin{thm}
\ollabel{thm:fseq-frm-equiv}
$\Frm[L]$ అనేది, నిర్మాణ క్రమం గల అన్ని $\Lang L$-సంకేతమాలలు~$!A$
ఉన్న సమితి.
\end{thm}

\begin{proof}
భాష~$\Lang L$లో నిర్మాణ క్రమం గల అన్ని సంకేతమాలల సమితిని $F$
అనుకుందాం. \olref[fol][syn][fseq]{prop:formed}లో
$\Frm[L] \subseteq F$ అని చూశాం; ఇప్పుడు తిరుగుదిశను నిరూపిస్తాం.

$!A$కు $\tuple{!A_0,\dotsc,!A_n}$ అనే నిర్మాణ క్రమం ఉందనుకుందాం.
$n$పై బలమైన ఆగమనం ద్వారా $!A \in \Frm[L]$ అని నిరూపిస్తాం. చివరి
సూచిక $m < n$ గల నిర్మాణ క్రమం ఉన్న ప్రతి సంకేతమాల $\Frm[L]$లో
ఉంటుందనేది మన ఆగమన పరికల్పన.
\sourcecorrection{OLTEFOLSYN-005}{ఆంగ్ల మూలం $n$పై ఆగమనం అంటూనే,
పరికల్పనను ``పొడవు $m<n$'' కోసం పేర్కొంది; కానీ
$\tuple{!A_0,\dotsc,!A_n}$ పొడవు $n+1$. పరికల్పనను చివరి సూచిక
$m<n$ అని స్పష్టపరచాం; ఇది వెంటనే వాడే $j,k<n$తో సరిపోతుంది.}
నిర్మాణ క్రమం నిర్వచనం ప్రకారం, $!A \ident !A_n$ పరమాణువు అయి ఉండాలి;
లేదా కింది వాటిలో ఒకటి నిజమయ్యేలా $j,k < n$లు ఉండాలి:
\begin{enumerate}
    \tagitem{prvNot}{$!A \ident \lnot !A_j$.}{}%
    \tagitem{prvAnd}{$!A \ident (!A_j \land !A_k)$.}{}%
    \tagitem{prvOr}{$!A \ident (!A_j \lor !A_k)$.}{}%
    \tagitem{prvIf}{$!A \ident (!A_j \lif !A_k)$.}{}%
    \tagitem{prvIff}{$!A \ident (!A_j \liff !A_k)$.}{}%
    \tagitem{prvAll}{$!A \ident \lforall[x][!A_j]$.}{}%
    \tagitem{prvEx}{$!A \ident \lexists[x][!A_j]$.}{}%
\end{enumerate}
ఇప్పుడు సందర్భాలవారీగా తర్కిద్దాం. $!A$ పరమాణువు అయితే
$!A_n \in \Frm[L]$.
\sourcecorrection{OLTEFOLSYN-004}{సిద్ధాంతమూ నిరూపణమూ భాష~$\Lang L$ను
నిర్దేశించినా, ఆంగ్ల మూలం పరమాణు సందర్భంలోనూ తరువాతి ద్విస్థానిక
సందర్భంలోనూ $\Frm[L_0]$కు మారింది. రెండింటినీ సందర్భం కోరే
$\Frm[L]$గా చేశాం.}
బదులుగా $!A \ident (!A_j \land !A_k)$ అనుకుందాం.
\sourcecorrection{OLTEFOLSYN-006}{ఆంగ్ల మూలం ఇక్కడ అర్థపర తుల్యత
$\equiv$ను వాడింది. నిర్మాణ క్రమ నిర్వచనం సంకేతమాల తాదాత్మ్యం
$\ident$ను కోరుతుంది కనుక, అదే సంబంధాన్ని నిలకడగా వాడాం.}
\olref[fol][syn][fseq]{lem:fseq-init} ప్రకారం,
$\tuple{!A_0,\dotsc,!A_j}$, $\tuple{!A_0,\dotsc,!A_k}$లు వరుసగా
$!A_j$,~$!A_k$లకు నిర్మాణ క్రమాలు. ఇవి $!A$ నిర్మాణ క్రమానికి నిజ
ఆద్య ఉపక్రమాలు కనుక వాటి చివరి సూచికలు~$n$కంటే తక్కువ. అందువల్ల
ఆగమన పరికల్పన ప్రకారం $!A_j$,~$!A_k$లు $\Frm[L]$లో ఉంటాయి; \tetoken{ఒక సూత్రం}{!!a{formula}}
నిర్వచనం ప్రకారం $(!A_j \land !A_k)$ కూడా అందులోనే ఉంటుంది. మిగతా
సందర్భాలు సమాంతర తర్కంతో వస్తాయి.
\end{proof}

పదాల నిర్మాణ క్రమాలకు కూడా \tetoken{సూత్రాలు}{!!{formula}s} నిర్మాణ క్రమాల వంటి ధర్మాలు
ఉంటాయి.

\begin{prop}
\ollabel{prop:fseq-trm-equiv}
$\Trm[L]$ అనేది పద నిర్మాణ క్రమం గల అన్ని $\Lang L$-సంకేతమాలలు~$t$
ఉన్న సమితి.
\end{prop}

\begin{proof}
సాధన.
\end{proof}

\begin{prob}
\olref[fol][syn][fseq]{prop:fseq-trm-equiv}ను నిరూపించండి.
సూచన: \olref[fol][syn][fseq]{thm:fseq-frm-equiv} నిరూపణలో వాడిన
వ్యూహాన్ని పోలిన వ్యూహాన్ని వాడండి.
\end{prob}

నిర్మాణ క్రమాల్లో రెండు రకాల ``వ్యర్థ'' భాగాలు ఉండవచ్చు: పునరావృత
మూలకాలు; సూత్రం లేదా పదం నిర్మాణానికి సంబంధం లేని మూలకాలు. కనిష్ఠ
నిర్మాణ క్రమాలను పరిశీలించడం ద్వారా రెండింటినీ తొలగించవచ్చు.

\begin{defn}[కనిష్ఠ నిర్మాణ క్రమాలు]
\ollabel{defn:minimal-fseq}
సూత్రం~$!A$కు నిర్మాణ క్రమం $\tuple{!A_0, \ldots, !A_n}$,
$!A$కు \emph{కనిష్ఠ నిర్మాణ క్రమం} కావడం అంటే $!A$కు ఉన్న ప్రతి ఇతర
నిర్మాణ క్రమం~$s$ కోసం, $!A$కు ఆ క్రమం~$s$ పొడవు $n+1$కంటే ఎక్కువ
లేదా దానికి సమానం కావడం.

అదే విధంగా, పదం~$t$కు నిర్మాణ క్రమం $\tuple{t_0, \ldots, t_n}$,
$t$కు \emph{కనిష్ఠ నిర్మాణ క్రమం} కావడం అంటే $t$కు ఉన్న ప్రతి ఇతర
నిర్మాణ క్రమం~$s$ కోసం, $t$కు ఆ క్రమం~$s$ పొడవు $n+1$కంటే ఎక్కువ లేదా
దానికి సమానం కావడం.
\end{defn}

ఒక సూత్రం లేదా పదానికి ఒకటికంటే ఎక్కువ కనిష్ఠ నిర్మాణ క్రమాలు
ఉండవచ్చని గమనించండి; కానీ వాటిలో సరిగ్గా అవే సంకేతమాలలు ఉంటాయి.

\begin{prop}
\ollabel{prop:subformula-equivs}
కిందివి తుల్యమైనవి:
\begin{enumerate}
    \item $!B$, $!A$కు ఉప-\tetoken{సూత్రం}{!!{formula}}.
    \item $!A$కు ఉన్న ప్రతి నిర్మాణ క్రమంలో $!B$ వస్తుంది.
    \item $!A$కు ఉన్న ఒక కనిష్ఠ నిర్మాణ క్రమంలో $!B$ వస్తుంది.
\end{enumerate}
\end{prop}

\begin{proof}
సాధన.
\end{proof}

\begin{prob}
\olref[fol][syn][fseq]{prop:subformula-equivs}ను నిరూపించండి.
\end{prob}

\begin{history}
రేమండ్ స్మల్యన్ తన \emph{First-Order Logic} పాఠ్యగ్రంథంలో
\citep{Smullyan1968} నిర్మాణ క్రమాలను ప్రవేశపెట్టారు. నిర్మాణ క్రమాల
అదనపు ధర్మాలను \citet{Zuckerman1973} స్థాపించారు.
\end{history}

\end{document}
